\chapter{Introduction}
\label{chap:introduction}

\section{Background and Motivation}
Civil and structural bridge design demands stringent adherence to safety standards, serviceability criteria, and numerical reliability. In India, bridge design is governed by the standards formulated by the Indian Roads Congress (IRC) and the Bureau of Indian Standards (BIS). 

The \textbf{OsdagBridge} project, part of the broader Osdag ecosystem supported by the FOSSEE initiative at IIT Bombay, represents a dedicated open-source platform tailored to the automated design, grillage analysis, 3D CAD modeling, and drafting of steel and composite bridge superstructures.

\section{The FOSSEE Screening Task}
As part of the FOSSEE Winter Internship 2026 screening task, this initiative focuses on rigorous software quality assurance, boundary condition verification, numerical stability evaluation, and defect investigation. The candidate repository (\url{https://github.com/Pranshu10khtk/OsdagBridge}) targets the upstream repository (\url{https://github.com/Nidhikhare12/OsdagBridge}).

\section{Report Objectives and Structure}
This technical report serves to:
\begin{enumerate}
    \item Detail the architecture, mathematical models, and design standards implemented within OsdagBridge.
    \item Establish a robust testing workspace and verification methodology covering basic geometry, custom material definitions, additional loading parameters, multi-run state isolation, dictionary data integrity, and CAD/plot visual output.
    \item Provide an uncompromised, objective testing framework that captures actual system behaviors without fabrication or artificial synthesis.
    \item Diagnose identified software anomalies down to the source file and line number, offering actionable, code-level fix recommendations.
\end{enumerate}
